\section{Conclusiones Generales}

Este trabajo se propuso investigar la viabilidad de los modelos de lenguaje de gran escala (LLMs) de código abierto como herramienta de apoyo para docentes de inglés como segunda lengua, abordando dos tareas centrales: la clasificación automática de textos según niveles del CEFR y la adaptación de textos hacia niveles de menor complejidad. A lo largo de los capítulos anteriores se presentaron resultados experimentales que permiten extraer conclusiones de alcance tanto técnico como pedagógico.

\subsection*{Sobre la clasificación automática de niveles CEFR}

Los experimentos de clasificación demostraron que los LLMs generativos de código abierto constituyen una alternativa robusta para la asignación automática de niveles CEFR, siempre que se implemente una estrategia de ajuste fino adecuada. El modelo Meta-Llama-3.1-8B-Instruct ajustado logró superar a herramientas comerciales especializadas como Text Inspector y alcanzar un rendimiento comparable al del clasificador de TSAR, cuyo desempeño sobresaliente en el conjunto de holdout se explica, con alta probabilidad, por un conocimiento previo de los datos de Cambridge incorporado durante su entrenamiento.

Un hallazgo transversal a todos los modelos evaluados fue la dificultad sistemática para distinguir niveles adyacentes, en especial en las fronteras B1--B2 y B2--C1. Esta limitación refleja no solo una restricción técnica de los modelos, sino también una ambigüedad intrínseca al propio marco CEFR, cuyos descriptores cualitativos no establecen fronteras nítidas entre niveles contiguos. En consecuencia, aun cuando los modelos exhiben un sesgo hacia niveles intermedios, predominantemente B2, este comportamiento no resulta pedagógicamente neutro: entregar a un estudiante un texto un nivel por encima del suyo puede tener valor formativo, mientras que una discrepancia de dos o más niveles resulta inadecuada.

Las estrategias de ingeniería de prompts mostraron una utilidad moderada y variable según el modelo. Si bien la incorporación de descriptores CEFR, listas de vocabulario graduado y ejemplos seleccionados por docentes mejoró el rendimiento en ciertos casos, esta mejora nunca alcanzó la magnitud obtenida mediante fine-tuning. Más aún, en modelos de menor escala o con ventanas de contexto más reducidas, la extensión excesiva del prompt derivó en una degradación del desempeño, posiblemente por efectos de saturación de atención. Este resultado matiza la expectativa de que más información siempre conduce a mejores resultados, y subraya la importancia de calibrar cuidadosamente tanto la cantidad como la calidad de la información provista al modelo.

\subsection*{Sobre la adaptación controlada por nivel de legibilidad}

Los experimentos de adaptación, enmarcados en la shared task TSAR 2025, confirmaron que los LLMs de parámetros reducidos son capaces de producir adaptaciones pedagógicamente aceptables cuando se los dota de recursos lingüísticos explícitos y estrategias de prompting elaboradas. El sistema \textit{TaskGen}, basado exclusivamente en el modelo Llama-3.1-8B-Instruct, logró posicionarse en el 6º lugar del ranking general por equipos, compitiendo con sistemas que emplearon modelos privativos de última generación o modelos abiertos de más de 70 mil millones de parámetros. Este resultado tiene implicancias significativas para contextos de investigación con recursos computacionales limitados, como el entorno universitario argentino en el que se desarrolló este trabajo.

La estrategia de \textit{ensemble} por votación heurística (Run~3) resultó ser la más robusta de las tres evaluadas, al lograr el mejor equilibrio entre la fidelidad semántica respecto al texto original y la similitud con las adaptaciones de referencia producidas por expertos. Este resultado valida la intuición de que ningún prompt único captura de manera óptima todos los aspectos relevantes de la tarea, y que la diversificación de estrategias, combinada con un mecanismo de selección bien calibrado, puede suplir en parte las limitaciones de modelos de menor escala.

\subsection*{Contribuciones del trabajo}

Más allá de los resultados experimentales, este trabajo realiza contribuciones concretas en tres planos. En primer lugar, en el plano metodológico, se propone y valida una métrica de evaluación propia, la Accuracy Aproximada, que incorpora la naturaleza ordinal del CEFR y permite una valoración pedagógicamente más informativa del error de clasificación, complementando al RMSE utilizado en el dominio de la adaptación. En segundo lugar, en el plano de los recursos, se compiló y puso a disposición pública un conjunto de recursos lingüísticos graduados por nivel CEFR, descriptores, listas de vocabulario y ejemplos morfosintácticos, que resultaron útiles tanto para los experimentos de clasificación como de adaptación, y que constituyen un insumo reutilizable para la comunidad investigadora. En tercer lugar, en el plano de los datos, se inició la construcción de un corpus propio de adaptaciones alineadas al CEFR, elaborado en colaboración con docentes expertos, que representa un recurso inédito en el dominio y cuya expansión futura podría transformar el estado del arte en la tarea de adaptación controlada para el inglés como L2.

\subsection*{Reflexión final}

Los resultados obtenidos demuestran que la brecha entre los sistemas de código abierto de tamaño moderado y los modelos privativos de gran escala, si bien existe, es menos amplia de lo que podría anticiparse cuando se invierte esfuerzo en el diseño de prompts, la selección de datos de entrenamiento y la evaluación rigurosa. En el contexto de la enseñanza de lenguas, donde la disponibilidad de herramientas accesibles, auditables y adaptables resulta tan importante como su precisión, esta constatación abre un camino promisorio para el desarrollo de recursos pedagógicos automatizados que puedan desplegarse sin depender de infraestructuras comerciales cerradas.

Al mismo tiempo, los desafíos identificados, la dificultad para discriminar niveles adyacentes, la tensión entre adecuación lingüística y preservación semántica, y la escasez de corpus paralelos alineados al CEFR, revelan que queda un camino considerable por recorrer antes de que estos sistemas puedan considerarse herramientas de apoyo docente plenamente confiables. La continuidad de esta línea de investigación, alimentada por corpus más ricos, estrategias de fine-tuning más expresivas y una colaboración sostenida con expertos pedagógicos, constituye la ruta más prometedora para cerrar esa brecha.